\subsection{Native validation and paired measurements}

The pinned baseline is \code{aeb891e7159a}. The driver
\code{weighted\_research/compact\_coordinates.py} records 442 current source
and fixture hashes, separately records every baseline package hash, and checks
that the source hashes do not change during a run. The dense baseline is
loaded in an isolated package namespace from Git; it includes the same
adaptive interval-merger implementation. The only changed pre-existing
runtime modules are the component geometry, producer and verifier; the
compact producer decoder is a new module.

The frozen Regina 7.4 corpus contains 1,275 vectors over 45 triangulations.
All three census modes and the normalized disc-count query pass 5,100 oracle
comparisons and independent certificate replays. The extra compact audit
checks full-summary equality against the dense predecessor on every vector,
replays all 1,275 dense certificates with the current verifier, and confirms
that every unweighted orbit proof is unchanged. The source-support type
bounds pass on every corpus entry, including 252 inputs with nonorientable
components. These are finite audit checks accompanying the proof above.
The audit takes 35.548 seconds, including the additional predecessor runs;
it overlaps the test suite and is not a performance measurement.

The maintained suite passes all 1,125 tests in 220.141 seconds. After tests,
audit, commit and publication finished, the benchmark ran serially for
267.924 seconds. It uses six supplied surfaces, four shuffled arms (dense,
compact, and an A/A copy of each), five measured rounds and one retained
warm-up round. All 120 measured calls and 24 warm-ups complete successfully.
Every timed call includes fresh source validation, coordinate-weight
construction, orbit discovery, certificate production and independent replay.
Fixture creation, serialization and cross-arm equality checks are outside
the timers. Both arms return identical full component summaries. Raw samples,
source pins and deduplicated full certificates are retained in
\code{data/compact-coordinates-benchmark.json}.

\begin{center}\small
\begin{tabular}{lrrrrr}
\toprule
Input & Dense ms & Compact ms & Dense/compact & Dense A/A & Compact A/A \\
\midrule
meridian-1 & 1.003 & 0.961 & 1.035 & 0.986 & 1.009 \\
meridian-16 & 54.962 & 22.945 & 2.413 & 0.993 & 0.994 \\
meridian-64 & 2204.769 & 484.808 & 4.548 & 0.988 & 1.042 \\
meridian-128 & 15969.781 & 2744.039 & 5.851 & 0.991 & 0.993 \\
triangle-only & 5.332 & 4.005 & 1.331 & 0.988 & 0.993 \\
binary-1024 & 28.181 & 16.203 & 1.705 & 1.014 & 0.996 \\
\bottomrule
\end{tabular}
\end{center}


At 128 tetrahedra, dimension falls from 896 to 128, and the initial and peak
piecewise-constant weight-run count falls from 385 to 257. The underlying
133 orbit cycles and 21,008 replay events are unchanged. Median total time
falls from 15.970 seconds to 2.744 seconds; the median within-round ratio is
$5.851$, which need not equal the ratio of separate medians. Certificate
size falls from 2,504,362 to 1,867,476 bytes in the recorded compact JSON
encoding. This is a complete full-coordinate-query comparison, not a
comparison of different output modes.

The corresponding paired gains are $2.413$ at 16 tetrahedra and $4.548$ at
64. Dense A/A ratios range from $0.986$ to $1.014$ and compact A/A ratios
from $0.993$ to $1.042$. The one-tetrahedron $1.035$ ratio is small enough
that it should not be promoted as a meaningful gain. A triangle-only source
with two vertex-link types improves by $1.331$; an eight-tetrahedron meridian
scaled by $2^{1024}$ and supplemented with $2^{1024}+1$ vertex links improves
by $1.705$. These are six supplied-source cases, not a universal speedup
claim or new automatic unknot-recognition coverage.
